% !TeX program = tectonic
\documentclass[11pt]{article}
\IfFileExists{alius-guidelines-style.tex}{\usepackage[a4paper,top=27mm,bottom=24mm,left=25mm,right=25mm]{geometry}
\usepackage[T1]{fontenc}
\usepackage[utf8]{inputenc}
\usepackage[scaled=0.96]{helvet}
\renewcommand{\familydefault}{\sfdefault}
\usepackage{array}
\usepackage{enumitem}
\usepackage{fancyhdr}
\usepackage[hidelinks]{hyperref}
\usepackage{microtype}
\usepackage{needspace}
\usepackage{tabularx}
\usepackage{xcolor}
\usepackage{xurl}

\definecolor{aliusgreen}{HTML}{13823A}
\definecolor{aliusgray}{HTML}{7A7A7A}
\definecolor{aliusline}{HTML}{9A9A9A}

\hypersetup{
  colorlinks=true,
  linkcolor=aliusgreen,
  urlcolor=aliusgreen
}

\setlength{\parindent}{0pt}
\setlength{\parskip}{0.72em}
\setlength{\tabcolsep}{0pt}
\raggedbottom
\sloppy
\urlstyle{same}

\newcommand{\ALIUSFooterTitle}{ALIUS Bulletin Guidelines (2022)}

\pagestyle{fancy}
\fancyhf{}
\fancyfoot[L]{\footnotesize\color{aliusgray}\ALIUSFooterTitle}
\fancyfoot[R]{\footnotesize\color{aliusgray}\thepage}
\renewcommand{\headrulewidth}{0pt}
\renewcommand{\footrulewidth}{0.45pt}

\setlist[itemize]{
  leftmargin=1.35em,
  itemsep=0.18em,
  topsep=0.15em,
  parsep=0em
}
\setlist[enumerate]{
  leftmargin=2.05em,
  itemsep=0.42em,
  topsep=0.2em,
  parsep=0em
}
\setlist[description]{
  leftmargin=0em,
  labelsep=0.55em,
  itemsep=0.35em,
  topsep=0.2em,
  font=\normalfont\bfseries
}

\newcommand{\guideDivider}{%
  \par\noindent{\color{aliusline}\rule{\textwidth}{0.45pt}}\par
}

\newcommand{\guideMeta}[2]{%
  {\footnotesize\bfseries\MakeUppercase{#1}}\par
  {\footnotesize #2}\par\vspace{0.65em}
}

\newcommand{\guideHeader}[5]{%
  \thispagestyle{fancy}%
  \vspace*{1mm}
  \begin{center}
    {\Large\bfseries\MakeUppercase{ALIUS Bulletin}\par}
    \vspace{0.25em}
    {\normalsize *\par}
    \vspace{0.25em}
    {\normalsize\bfseries\MakeUppercase{Guidelines For}\par}
    \vspace{0.15em}
    {\LARGE\bfseries\MakeUppercase{#1}\par}
  \end{center}
  \vspace{1.05em}
  \begin{tabularx}{\textwidth}{@{}>{\raggedright\arraybackslash}p{0.55\textwidth}!{\hspace{1.35em}{\color{aliusline}\vrule width 0.5pt}\hspace{1.35em}}>{\raggedright\arraybackslash}X@{}}
    {\large\itshape #2} &
    \guideMeta{Format}{#3}
    \guideMeta{Editorial Route}{#4}
    \guideMeta{Contact}{#5}
  \end{tabularx}
  \vspace{0.75em}
  \guideDivider
  \vspace{0.25em}
}

\newcommand{\guideSection}[1]{%
  \Needspace{6\baselineskip}%
  \par\vspace{1.1em}
  \begin{center}
    {\color{aliusline}\rule{0.14\textwidth}{0.45pt}}\hspace{0.85em}%
    {\large\bfseries #1}%
    \hspace{0.85em}{\color{aliusline}\rule{0.14\textwidth}{0.45pt}}
  \end{center}
  \vspace{-0.25em}
}

\newcommand{\guideSubsection}[1]{%
  \par\vspace{0.65em}{\bfseries #1}\par\vspace{-0.35em}
}

\newcommand{\guideQuestion}[1]{%
  \par\vspace{0.55em}{\color{aliusgreen}\bfseries #1}\par
}

\newenvironment{guideChecklist}
  {\begin{itemize}[label=\textendash]}
  {\end{itemize}}
}{\usepackage[a4paper,top=27mm,bottom=24mm,left=25mm,right=25mm]{geometry}
\usepackage[T1]{fontenc}
\usepackage[utf8]{inputenc}
\usepackage[scaled=0.96]{helvet}
\renewcommand{\familydefault}{\sfdefault}
\usepackage{array}
\usepackage{enumitem}
\usepackage{fancyhdr}
\usepackage[hidelinks]{hyperref}
\usepackage{microtype}
\usepackage{needspace}
\usepackage{tabularx}
\usepackage{xcolor}
\usepackage{xurl}

\definecolor{aliusgreen}{HTML}{13823A}
\definecolor{aliusgray}{HTML}{7A7A7A}
\definecolor{aliusline}{HTML}{9A9A9A}

\hypersetup{
  colorlinks=true,
  linkcolor=aliusgreen,
  urlcolor=aliusgreen
}

\setlength{\parindent}{0pt}
\setlength{\parskip}{0.72em}
\setlength{\tabcolsep}{0pt}
\raggedbottom
\sloppy
\urlstyle{same}

\newcommand{\ALIUSFooterTitle}{ALIUS Bulletin Guidelines (2022)}

\pagestyle{fancy}
\fancyhf{}
\fancyfoot[L]{\footnotesize\color{aliusgray}\ALIUSFooterTitle}
\fancyfoot[R]{\footnotesize\color{aliusgray}\thepage}
\renewcommand{\headrulewidth}{0pt}
\renewcommand{\footrulewidth}{0.45pt}

\setlist[itemize]{
  leftmargin=1.35em,
  itemsep=0.18em,
  topsep=0.15em,
  parsep=0em
}
\setlist[enumerate]{
  leftmargin=2.05em,
  itemsep=0.42em,
  topsep=0.2em,
  parsep=0em
}
\setlist[description]{
  leftmargin=0em,
  labelsep=0.55em,
  itemsep=0.35em,
  topsep=0.2em,
  font=\normalfont\bfseries
}

\newcommand{\guideDivider}{%
  \par\noindent{\color{aliusline}\rule{\textwidth}{0.45pt}}\par
}

\newcommand{\guideMeta}[2]{%
  {\footnotesize\bfseries\MakeUppercase{#1}}\par
  {\footnotesize #2}\par\vspace{0.65em}
}

\newcommand{\guideHeader}[5]{%
  \thispagestyle{fancy}%
  \vspace*{1mm}
  \begin{center}
    {\Large\bfseries\MakeUppercase{ALIUS Bulletin}\par}
    \vspace{0.25em}
    {\normalsize *\par}
    \vspace{0.25em}
    {\normalsize\bfseries\MakeUppercase{Guidelines For}\par}
    \vspace{0.15em}
    {\LARGE\bfseries\MakeUppercase{#1}\par}
  \end{center}
  \vspace{1.05em}
  \begin{tabularx}{\textwidth}{@{}>{\raggedright\arraybackslash}p{0.55\textwidth}!{\hspace{1.35em}{\color{aliusline}\vrule width 0.5pt}\hspace{1.35em}}>{\raggedright\arraybackslash}X@{}}
    {\large\itshape #2} &
    \guideMeta{Format}{#3}
    \guideMeta{Editorial Route}{#4}
    \guideMeta{Contact}{#5}
  \end{tabularx}
  \vspace{0.75em}
  \guideDivider
  \vspace{0.25em}
}

\newcommand{\guideSection}[1]{%
  \Needspace{6\baselineskip}%
  \par\vspace{1.1em}
  \begin{center}
    {\color{aliusline}\rule{0.14\textwidth}{0.45pt}}\hspace{0.85em}%
    {\large\bfseries #1}%
    \hspace{0.85em}{\color{aliusline}\rule{0.14\textwidth}{0.45pt}}
  \end{center}
  \vspace{-0.25em}
}

\newcommand{\guideSubsection}[1]{%
  \par\vspace{0.65em}{\bfseries #1}\par\vspace{-0.35em}
}

\newcommand{\guideQuestion}[1]{%
  \par\vspace{0.55em}{\color{aliusgreen}\bfseries #1}\par
}

\newenvironment{guideChecklist}
  {\begin{itemize}[label=\textendash]}
  {\end{itemize}}
}
\renewcommand{\ALIUSFooterTitle}{ALIUS Bulletin Guidelines for Reviews (2022)}

\begin{document}

\guideHeader
  {Reviews}
  {If you want to edit a review for the ALIUS Bulletin, please read the following guidelines carefully.}
  {Short review of a published work; 500--1,000 words plus references and author details.}
  {Reviewed by editors and proofreaders, then prepared for publication in the Bulletin.}
  {\href{mailto:aliusresearchgroup@gmail.com}{aliusresearchgroup@gmail.com}}

Reviews offer an outlook on published works on consciousness and their relevance with respect to the diversity of consciousness. A review is intended as a short commentary, composed of a summary of the main ideas of the work, 250--500 words, and an account of its relevance for the study of the diversity of consciousness, 250--500 words.

\guideSection{Length of the Review}

Each review must include the following:

\begin{guideChecklist}
\item A title.
\item A main text of 500--1,000 words.
\item References, with no limit on the number of references. The style used should be the 6th edition of APA. If you use Zotero, this citation style can be downloaded from the \href{https://www.zotero.org/styles?q=id%3Aapa}{Zotero APA style page}.
\item Addresses and affiliations of the authors for correspondence.
\item Font in the submitted file: Times New Roman, size 12.
\end{guideChecklist}

\guideSection{Content of the Review}

Reviews are intended to highlight how theories of consciousness, including recently formulated ones, fit with the diversity of consciousness. They should both highlight recent progress in consciousness science and provide a perspective on how these works can inform the study of the diversity of consciousness.

After describing one or several main ideas of the article or book, reviews can depart from the message of the original author and critically assess how the diversity of consciousness, or at least one specific conscious state, is given a new perspective in light of this theory. Reviews can also discuss unresolved theoretical points of the article or book from the perspective of the diversity of consciousness.

By doing so, reviews aim to raise the importance of embracing the diversity of consciousness when discussing the nature of cognition and/or consciousness. Each comment should be easily understandable to any scholar. For example, an anthropologist reading a comment conducted with a neuroscientist should easily understand most of the discussion, and conversely, a neuroscientist reading an interview conducted with an anthropologist should easily understand most of the discussion.

\guideSection{Modus Operandi and Requirements}

\begin{enumerate}
\item Ask one of the editors of the Bulletin whether the article or book you want to review is deemed relevant as an open review.
\item Once editors have received the review, they will forward it to a native English-speaking proofreader who will check grammar and linguistic correctness. Note that the proofreader's job is not to correct typos. Typos must be corrected before sending the file to the editors.
\item The editors will next send you the feedback provided by the proofreaders. You will be expected to integrate the changes suggested by the proofreaders.
\item Once you have the author or authors' agreement for publication, send the final version of the review to one of the editors. The review will then be ready to be published.
\end{enumerate}

\guideSection{Editorial Team}

ALIUS contact email: \href{mailto:aliusresearchgroup@gmail.com}{aliusresearchgroup@gmail.com}.

Daniel Friedman (\href{mailto:danielarifriedman@gmail.com}{danielarifriedman@gmail.com}); Matthieu Koroma (\href{mailto:matthieu.koroma@gmail.com}{matthieu.koroma@gmail.com}); George Fejer (\href{mailto:gyorgyf@live.com}{gyorgyf@live.com}).

\end{document}
